\section{Application Overview}
\label{sec:app_overview}

The application is a multimodal book discovery platform. Users can search by typing a query, uploading a cover image, or combining both, and receive personalised recommendations that adapt as they interact with the catalogue.

Three design decisions shaped what the application needs to do well.
The catalogue covers roughly 3 million items, so unguided browsing is not practical and all discovery paths go through embedding-based retrieval.
Vietnamese-speaking users can query the system directly; the NLLB-200 pipeline translates the input to English before encoding, with no change to the search interface.
Personalisation works from the first visit: a new user immediately receives co-purchase suggestions from the Cleora graph, and once five click interactions accumulate, DIF-SASRec sequential recommendations engage automatically.

The full application stack is:
\begin{itemize}
    \item Frontend: React 19 + Vite 8 + Tailwind CSS, served at \texttt{http://localhost:5173}.
    \item Backend: FastAPI + uvicorn, served at \texttt{http://localhost:8000}. Seven route groups: \texttt{/search}, \texttt{/recommend}, \texttt{/interact}, \texttt{/profile}, \texttt{/auth}, \texttt{/ask\_llm}, \texttt{/health}.
    \item Infrastructure: MongoDB (\texttt{nba\_logs} database) and Redis (interaction write-queue), both provisioned via Docker Compose.
\end{itemize}

If the backend is unreachable, a dismissible error notification appears and the interface continues showing the most recently loaded data.
User identity is persisted in \texttt{localStorage} so that returning users bypass the login gate automatically.
